% Build the layer, Chapter 8, second build: the transfer of a card.
% Inserted in \S\ref{sec:fh-xfer}, after \S\ref{sec:fh-xfer-att}, which
% states the bound on what a transfer may grant. This section builds the
% four-message ceremony in the kernel's decision core, for the case of
% \S\ref{sec:fh-ex-task}, and runs it.
%
% No code or output in this file is typed by hand. Each listing is a file
% under excerpts/, which the kernel's tools/excerpts.mjs cuts from its
% source and its recorded runs (see figures/pd-listing.tex). To change a
% listing, change the kernel and cut again. The prose refers to functions
% by name and never by line number, so a listing that moves within its file
% does not make a sentence wrong.
%
% \pdlistingroom{n} before a listing of n lines keeps it on one page. When a
% listing changes length, change the number with it.
\subsection{Build the layer: the transfer of a card}
\label{sec:fh-offer}

\noindent\emph{Express lane: Heron holds a card from Alice's harbor, and
its migration is at Carol's. A gate that honors Alice's cards as they are
admits whatever Alice's harbor signs. In the build, Alice's harbor signs an
offer that describes the card, and Carol's harbor issues a card of its own
for what its own mapping grants. It honors an offer once. It revokes its
card when Alice's harbor says, with its signature, that the old card is
revoked. An exhaustive test of two harbors reaches 195 states and finds no
card that should not be there. In 66 of them Alice's harbor has revoked the
card, and a harbor that has not been told still admits its own.}

Section~\ref{sec:fh-xfer-ceremony} gives the ceremony as four messages, and
nothing in the chapter has run them. This section builds the ceremony in
\pdkernelname{} for the task of \S\ref{sec:fh-ex-task}, and runs it.

The cast is that task's. Heron is one of Alice's workers. Its card $C_A$
lets it write the \texttt{checkout} schema of Alice's test database, which
the runs write as \texttt{test-db/checkout/*}. From
Listing~\ref{alg:fh-offer-run-core} on, the runs call that card
\texttt{alice-1}. The migration is on Carol's staging database,
\texttt{staging-db}. Finch is another of Alice's workers. Bob's harbor takes offers from Alice's as well. Wren is an
outsider with a copy of every message.

The section builds the offer and the card that answers it. It also builds
the revocation that follows the card. It leaves out the epoch root that
the chapter's offer carries, and the boundary at the end lists what else
is missing.

\subsubsection{Naive: honor the other harbor's cards}

Carol's gate admits cards that Carol's harbor signed. The shortest way to
let Heron in is to give the gate a second key, so that it honors a card
from Alice's harbor as it honors its own
(Listing~\ref{alg:fh-offer-honor}).

\pdlistingroom{9}
\pdexcerpt[caption={The first attempt. A gate honors another harbor's
key.},label={alg:fh-offer-honor}]{ch08b-honor}

Heron's card names tables of Alice's test database, and those are no use
at Carol's. So Alice's harbor signs Heron a second card, for
\texttt{staging-db/checkout/*}.

\subsubsection{Break: a card for the whole database}

Heron's migration goes through. Then Alice's harbor signs Finch a card
for every table of the staging database
(Listing~\ref{alg:fh-offer-run-foreign}).

\pdlistingroom{17}
\pdrun[caption={A card that only Alice's harbor signed, at Carol's gate.
The word at the left of each step says who acts, and \texttt{alice} and
\texttt{carol} are the two harbors.},label={alg:fh-offer-run-foreign}]{transfer-foreign}

Nothing was forged, and Carol's gate did what it was built to do. At
step~5 it admits a write to a billing table that is no part of the job.
Whatever Alice's harbor writes in a card is now authority at Carol's, and
Carol's harbor decided none of it. The example of
\S\ref{sec:fh-sovereignty} draws this line: who signs a card is one
question, and who rules on what the card touches is another.

\subsubsection{Fix: an offer, and a card of the receiver's own}

In the ceremony the card never leaves Alice's harbor. That harbor signs an
offer instead (Listing~\ref{alg:fh-offer-text}). An offer names both
harbors by their keys. It names the card by its id and the holder by its
key. It lists the rights offered, and says when the card expires. It has
an identifier, \texttt{xid}, and it names the mapping it asks under. The
signature covers every one of those fields.

What an offer is worth is for the harbor that receives it to say. Carol's
harbor keeps a mapping for Alice's, which is a list of rules. A rule has a
right that Alice's harbor may offer, and the right that Carol's harbor
grants for it. Carol's one rule turns a write to
\texttt{test-db/checkout/*} into a write to
\texttt{staging-db/checkout/*}.

\pdlistingroom{20}
\pdexcerpt[caption={The text an offer is a signature over, and the check
of the signature.},label={alg:fh-offer-text}]{ch08b-offer}

\texttt{granted} applies the rules (Listing~\ref{alg:fh-offer-granted}).

\pdlistingroom{13}
\pdexcerpt[caption={What a harbor grants for an
offer.},label={alg:fh-offer-granted}]{ch08b-granted}

\texttt{granted} never copies a right out of the offer. This is the bound
of Design invariant~\ref{inv:fh-transfer}: the new card holds what the
mapping makes of the offer, and no more. \texttt{accept} puts the checks
in order (Listing~\ref{alg:fh-offer-accept}).

\pdlistingroom{34}
\pdexcerpt[caption={A harbor is sent an
offer.},label={alg:fh-offer-accept}]{ch08b-accept}

\texttt{accept} checks that the offer is signed by the harbor it names as
its sender, and that it names this harbor as its recipient. The harbor
must have a mapping for the sender, and the offer must name the version of
it that the harbor has now. The offered card must not have expired. The
checks \texttt{used} and \texttt{revoked} belong to the two breaks below.

The core holds no key
(\weld{prereq:sec:prereq-build}{decision core}), so the harbor issues the
card (Listing~\ref{alg:fh-offer-receive}).

\pdlistingroom{25}
\pdexcerpt[caption={A harbor asks the core about an offer, and issues the
card.},label={alg:fh-offer-receive}]{ch08b-receive}

\texttt{receive} gives the core this harbor's key and its mapping for the
sender, with two facts it has looked up. It signs a card
with the claims the core returns, and records which offer the card came
from. The card is an ordinary card of Carol's harbor. Her gate checks it
as it checks any other, and Heron has to prove it holds the card's key
(\weld{anchor:sec:anchor-build-holder}{proof of possession}).

Listing~\ref{alg:fh-offer-run-core} is the run. Steps 4 to 8 carry the
four messages. At step~4 Heron asks its own harbor, and proves it holds
the card by answering the gate's challenge. Step~5 is the offer. At
step~6 Carol's harbor issues \texttt{carol-1}, and at step~7 it tells
Alice's harbor which card it issued. The card goes to Heron at step~8.
Carol's gate refused Heron's old card at step~3, and it admits the new
one at step~9.

The run goes on to two wider cards. Finch's card for the whole test
database buys what Heron's did, because the mapping has one rule. The
card that broke the first attempt is offered at step~14 and buys nothing,
since no rule of Carol's has \texttt{staging-db} on its offered side.
Alice's harbor can still sign whatever it likes. Its signature is no
longer authority at Carol's.

\pdlistingroom{39}
\pdrunlines[caption={The four messages of a transfer, and then two wider
cards, at harbors that ask the
core.},label={alg:fh-offer-run-core}]{transfer-core}{1,5-42}{1 and 5\textendash 42}

\subsubsection{Break again: the same offer, twice}

An offer is a message, and a message can be copied. The harbor above, as
first written, kept no record of the offers it had taken. Carol's harbor
revokes the card it issued, and Wren sends the offer again
(Listing~\ref{alg:fh-offer-run-replay}).

\pdlistingroom{20}
\pdrunlines[caption={An offer arrives a second time. Steps 1 and 2, left
out, are those of
Listing~\ref{alg:fh-offer-run-core}.},label={alg:fh-offer-run-replay}]{replayed-offer-offers}{1,9-27}{1 and 9\textendash 27}

At step~7 the harbor issues \texttt{carol-2}, and Heron is back in. Wren
gained nothing for herself, because the card names Heron's key. The loss
is Carol's. Her harbor revoked a card, and anyone with a copy of the offer
can undo that.

\subsubsection{Fix: honor an identifier once}

The harbor now keeps what it issued for each identifier.
\texttt{receive} looks the identifier up and passes the answer to the
core as \texttt{used}, and \texttt{accept} answers a used one with
\texttt{Known} (Listing~\ref{alg:fh-offer-run-once}).

\pdlistingroom{10}
\pdrunlines[caption={The same run at a harbor that honors an identifier
once. The steps left out are the first six of
Listing~\ref{alg:fh-offer-run-replay}.},label={alg:fh-offer-run-once}]{replayed-offer-core}{1,19-27}{1 and 19\textendash 27}

\texttt{Known} is an answer, and not a refusal. The harbor sends back the
acknowledgment of the card it issued the first time. A sender whose
acknowledgment was lost can send its offer again and learn what happened.
The log of \S\ref{sec:fh-log} answers a root it already holds in the same
way.

\subsubsection{Break again: revoked at home}

Alice revokes Heron's card, as she does at the end of the worked example
(\S\ref{sec:fh-ex-revocation}). Carol's harbor issued its card before
that, and nothing ties the two together
(Listing~\ref{alg:fh-offer-run-home}).

\pdlistingroom{29}
\pdrunlines[caption={Alice's harbor revokes the card that Carol's harbor
issued from. Steps 1 and 2 are left out
again.},label={alg:fh-offer-run-home}]{revoked-at-home-once}{1,9-36}{1 and 9\textendash 36}

Heron is refused at Alice's gate at step~6 and admitted at Carol's at
step~7. Step~8 changes nothing, because this harbor keeps no record of
another harbor's revocations. Step~10 is worse. A second offer of the
card, signed before the revocation, arrives after it, and the harbor
issues a card from a card that is revoked.

\subsubsection{Fix: follow the card's own harbor}

Carol's harbor has a record of which card each of its cards came from.
When it is told that one of those is revoked, it checks who is telling it
(Listing~\ref{alg:fh-offer-parent}).

\pdlistingroom{21}
\pdexcerpt[caption={A harbor is sent a revocation of another harbor's
card.},label={alg:fh-offer-parent}]{ch08b-parent}

The signed text is the one that Alice's own gate checks when a card is
revoked there
(\weld{anchor:sec:anchor-build-warranted}{signed revocation}), so one
signature serves at both harbors. The two differ in whose signature they
take. Alice's gate takes her harbor's, and also that of a holder the
card's chain names. Carol's harbor never saw the chain. It knows one key
that can speak for the card, and that is the key of the harbor that
offered it.

\texttt{told} is what the harbor does with a revocation that passes
(Listing~\ref{alg:fh-offer-told}). It does two things. It revokes every
card this harbor issued from the old one. It also remembers the old card,
and \texttt{accept} refuses any later offer of it. The second is for the
offer that was held up (Listing~\ref{alg:fh-offer-run-told}).

\pdlistingroom{20}
\pdexcerpt[caption={What a harbor does with a revocation that
passes.},label={alg:fh-offer-told}]{ch08b-told}

\pdlistingroom{14}
\pdrunlines[caption={The same run at harbors that ask the core. The steps
left out are the first seven of
Listing~\ref{alg:fh-offer-run-home}.},label={alg:fh-offer-run-told}]{revoked-at-home-core}{1,21-33}{1 and 21\textendash 33}

Step~9 is now a refusal, and step~10 refuses the late offer. Step~7 is as
it was. Between Alice's revocation and its arrival, Carol's gate admitted
Heron once. The build shortens that window to the time one message takes,
and does not close it. Theorem~\ref{thm:fh-conv} is about how long that
time is.

\subsubsection{Verify: an outsider, then every order}

\paragraph{An outsider.}
Wren has a copy of every message, the offer and the new card among them
(Listing~\ref{alg:fh-offer-run-outsider}).

\pdlistingroom{34}
\pdrunlines[caption={An outsider with a copy of every message. Steps 1 to
4, left out, set up the transfer and give Bob's harbor a
mapping.},label={alg:fh-offer-run-outsider}]{false-offer-core}{1,14-46}{1 and 14\textendash 46}

At step~2 Bob's harbor took a mapping for Alice's too, with the same name
and version as Carol's. So at step~5 the recipient's name is all that
stands between Wren and a card at Bob's, and \texttt{addressed} is the
check that reads it. Step~6 is an offer that Wren signed, and step~7 is
Alice's offer with one field changed. Neither has a signature by Alice's
harbor over what it says. Step~8 is \texttt{Known}. At step~9 she
presents Heron's card and cannot answer for Heron's key. Steps 10 and 11
are a revocation and an acknowledgment that she signed herself.

\paragraph{Every order.}
The order in which offers and revocations arrive is the network's to
choose, or an attacker's, so a few runs are not enough. The exhaustive
test sends two harbors seven offers and two revocations of
\texttt{alice-1} in every order, and the heading of
Listing~\ref{alg:fh-offer-walk} says which. In every state the test
compares each harbor with a small reference model. The model checks no
signature. It is told which offer was made for which harbor, and it keeps
its own account of what each harbor has been sent.

\pdlistingroom{37}
\pdrun[caption={Every sequence of events for two harbors, at each design
that takes offers.},label={alg:fh-offer-walk}]{transfers}

Three of the offers are good: two of Heron's card to Carol's harbor, and
one of it to Bob's. The other four should buy nothing. A harbor that
issues a card for every offer would issue without end, so the test stops
sending an offer to a harbor that has issued two cards for it. Even so,
that design holds a card it should not in 7,455 of its 7,862 states.
Honoring an identifier once brings that column to zero. It leaves
\texttt{revived}: in 224 of 407 states a gate admits a card after its
harbor was told that the old one is revoked. The harbors that ask the
core have neither fault in any of their 195 states.

The column \texttt{behind} is not zero for any design, and it cannot be.
In 66 of the core's states Alice's harbor has revoked the card, and a
harbor that has not been told admits its own. The last column says what
one notice does. At the core's harbors it ends all 66. \emph{Now you
try:} Exercise~\ref{ex:fh-offer-late} follows the offer that was held up
through \texttt{told} and \texttt{accept}.

\paragraph{Each guard removed.}
The build has nine guards. Eight of them are in the path of this test,
and the core can be compiled with any one removed
(Listing~\ref{alg:fh-offer-guards}). The ninth is the signature on an
acknowledgment. The test sends no acknowledgments, and unit tests in the
kernel cover that guard.

\pdlistingroom{11}
\pdrunlines[caption={The exhaustive test on the core, and on eight builds
of it with one guard removed from each. The heading and the notes on the
columns, left out, are those of
Listing~\ref{alg:fh-offer-walk}.},label={alg:fh-offer-guards}]{transfers-guards}{1,15-24}{1 and 15\textendash 24}

Every row differs from the first. A core that takes an offer whichever
harbor it names reaches 19,491 states, and in nearly all of them a harbor
holds a card for an offer that was made for the other one. A core that
grants the rights an offer names issues cards that Carol's gate then
refuses for the migration, which is the 601 under \texttt{refused}. The
row \texttt{late-offer} is the offer that was held up: without the memory
of revoked cards, a gate admits a card in 130 states whose parent the
harbor knows is revoked. A core that takes a revocation from any key lets
Wren cost Heron its card in 382 states.

\paragraph{The offer, as a protocol.}
The exhaustive test uses seven offers that the kernel made. ProVerif asks
about any offers an attacker can make
(\weld{anchor:sec:correct}{ProVerif}). In the model the attacker chooses
what each offer says and which harbor it is for, and Alice's harbor signs
it as often as it is asked. The attacker sees every message and has a
key of its own (Listing~\ref{alg:fh-offer-models}).

\pdlistingroom{29}
\pdrun[caption={ProVerif~2.05 on an offer and its revocation, and on five
controls.},label={alg:fh-offer-models}]{models-transfer}

\texttt{offer.pv} is the check as the core makes it. Its first result is
an injective correspondence. Each card that a harbor issues answers an
offer that Alice's harbor made for that harbor, and no two cards answer
one offer. Its second is \texttt{false} on purpose: a card can be issued,
so the first result is not true merely because nothing happens. Its third
says that a harbor acts on a revocation only if Alice's harbor signed it.

``No two cards'' rests on the rule that a harbor honors an identifier
once. ProVerif has no table that is read and written in one step, so the
model states the rule as a restriction on what a harbor's process may do.
\texttt{reuse-offer.pv} is the model without it. ProVerif finds the
replay there, and prints a second line: the same query without
\texttt{inj-} still holds. Every card still answers an offer that was
made, and two cards answer the same one.

\keyidea[A card of your own, for an offer you checked]{A harbor does not
honor another harbor's cards. It takes an offer, and decides from its own
mapping what the offer is worth. The card it issues is its own to revoke.
The offer is good once, at the harbor it names. The new card expires when
the old one does, and is revoked when the harbor learns that the old one
is.}

\pdscope{The exhaustive test covers two harbors and the seven offers it
names, at one moment. It counts states and says nothing about how likely
a state is, or how long a revocation takes to arrive. The ProVerif models
cover an offer and a revocation. They have no mapping and no expiry, and
they say nothing about the order in which an offer and a revocation
arrive. They assume that an identifier is honored once, and the
exhaustive test is what checks that the kernel does so.}

\begin{pdboundary}
\setlength{\parskip}{4pt plus 1pt}
The offer carries no epoch root. The chapter's offer names the root under
which the card was issued, and the receiving harbor checks that root
against a witness log. That needs the proof that an event is under a
root, which \S\ref{sec:fh-log} does not build. So Carol's harbor takes
Alice's harbor at its word that the card exists.

Heron signs its request, and what it signs does not name Carol's harbor.
Whoever carries the request to Alice's harbor can name Bob's instead
(Exercise~\ref{ex:fh-offer-request}).

A revocation reaches another harbor when someone sends it. These runs
send it by hand. Section~\ref{sec:fh-revocation} is about what gossip can
promise, and what it cannot.

The chapter's card carries the identifier of the offer it answers. The
card built here is an ordinary card of its harbor, and the harbor keeps
the record of which offer it came from. A gate that has only the card
cannot tell.

A harbor never forgets an identifier, though it could drop one once its
card has expired. The acknowledgment is checked, and no later step
depends on it. No bond is posted for a transfer and nothing is paid,
which is the subject of \S\ref{sec:fh-settlement}.

The two harbors run in one process, though the logs of
\S\ref{sec:fh-log} can each run as a process of its own.
\end{pdboundary}
